\documentclass{article}
\usepackage{amsmath,amssymb}
\usepackage{xurl}
\usepackage[hidelinks]{hyperref}
% Shared commands for notes in docs/paper-gaps/.

\DeclareUnicodeCharacter{2080}{\ifmmode{}_{0}\else\textsubscript{0}\fi}
\DeclareUnicodeCharacter{2081}{\ifmmode{}_{1}\else\textsubscript{1}\fi}
\DeclareUnicodeCharacter{2082}{\ifmmode{}_{2}\else\textsubscript{2}\fi}
\DeclareUnicodeCharacter{2099}{\ifmmode{}_{n}\else\textsubscript{n}\fi}

\newcommand{\Lean}{\textsc{Lean}}
\ExplSyntaxOn
\NewDocumentCommand{\leanid}{m}{
  \ifmmode
    \textsf{\detokenize{#1}}
  \else
    \group_begin:
    \urlstyle{sf}
    \tl_set:Nn \l_tmpa_tl {#1}
    \tl_replace_all:Nnn \l_tmpa_tl {\_} {_}
    \exp_args:NV \path \l_tmpa_tl
    \group_end:
  \fi
}
\ExplSyntaxOff
\providecommand{\tr}{\operatorname{tr}}
\newcommand{\tnleanIssueBase}{https://github.com/LionSR/TNLean/issues/}
\newcommand{\tnleanPrBase}{https://github.com/LionSR/TNLean/pull/}
\newcommand{\tnleanDocBase}{https://sirui-lu.com/TNLean/docs/}
\newcommand{\ghissue}[1]{\href{\tnleanIssueBase#1}{GitHub issue \##1}}
\newcommand{\ghpr}[1]{\href{\tnleanPrBase#1}{PR \##1}}
\newcommand{\doclink}[1]{\href{\tnleanDocBase#1.html}{\path{#1}}}

% Dirac notation and the identity operator, as in blueprint/src/macros/common.tex.
% \providecommand keeps note-local definitions authoritative.
\usepackage{dsfont}
\providecommand{\ket}[1]{|#1\rangle}
\providecommand{\bra}[1]{\langle #1|}
\providecommand{\braket}[2]{\langle #1 | #2 \rangle}
\providecommand{\ketbra}[2]{|#1\rangle\!\langle #2|}
\providecommand{\Id}{\mathds{1}}
\providecommand{\one}{\mathds{1}}
\providecommand{\C}{\mathbb{C}}
\providecommand{\MN}[1]{M_{#1}(\C)}
\providecommand{\MD}{M_D(\C)}

% Verdict marker: \gapnote{<kind>}{<status>}. Kind is the policy
% classification (clarification, local-correction, scope-restriction,
% unfaithful, false-source, open-gap); status is open, wip (elimination
% underway), resolved, or historical. Machine-read by the site index;
% prints a verdict line.
\newcommand{\gapnote}[2]{\par\noindent\textbf{Verdict:} #1 (#2).\par}

\title{From the Gates to the Tensors in the MPU Index Classification}
\author{}
\date{2026-11-16}
\begin{document}
\maketitle
\gapnote{clarification}{resolved}

\section{The printed argument}

The ``if'' half of clause (iv) of the Index Theorem of
Cirac--P\'erez-Garc\'ia--Schuch--Verstraete states that two matrix product
unitary tensors in canonical form with the same index are equivalent. Its
proof reads \cite[lines~857--865]{Cirac2017MPU}: after attaching ancillas and
blocking, both tensors are simple and have the same physical dimension; the
corresponding gates $u$ and $v$ then have the same input and output
dimensions, ``so that they can be smoothly connected to each other, and thus
so can'' the two blocked tensors.

The step from a path of gates to a path of tensors is not spelled out. A
path of gates joins the tensors constructed from the gates, but a given
tensor is not literally the tensor constructed from its own gates: the two
have different bond dimensions and different half factors, and only their
periodic operators agree.

\section{The supplied step}

Let $W$ be an MPU tensor in canonical form of bond dimension $D$, and let
$m\geq D^2$. The connection is made in three steps, all through tensors
generating MPUs.
\begin{enumerate}
    \item $W$ is similar, $W=X^{-1}YX$, to a tensor $Y$ in canonical form~II
          of the same bond dimension. A path $X(t)$ of invertible matrices from
          the identity to $X$, which exists because the singular points of
          $z\mapsto(1-z)I+zX$ are finitely many, joins $W_{2m}$ to $Y_{2m}$
          through the tensors $X(t)W_{2m}X(t)^{-1}$.
    \item $Y_m$ is simple, so $Y_{2m}$, the two-site block of $Y_m$, is in
          standard form with the gates $u,v$ of $Y_m$ and some half factors
          $(X_1,X_2)$. The tensor constructed from $u$ and $v$ is in standard
          form with the same gates and other half factors $(X_1',X_2')$.
    \item In the bond space of the sum of the two bond dimensions, the half
          factors $(\cos\theta\,X_1,\ \sin\theta\,X_1')$ and
          $(\cos\theta\,X_2,\ \sin\theta\,X_2')$ contract to
          $\cos^2\theta\,v+\sin^2\theta\,v=v$. Every intermediate tensor is
          therefore in standard form with gates $u,v$ and generates an MPU.
          For $\theta\in[0,\pi/2]$ this joins the two presentations, up to a
          permutation of bond coordinates, which is again a gauge.
\end{enumerate}
Equal indices and the product formula $r\ell=n^2$, with $n$ the physical
dimension of $Y_m$, force equal intermediate dimensions, and the paths of
unitary gates then complete the argument as printed.

\section{Formal statement}

The three steps are \leanid{MPOTensor.IsMPU.joinedIn\_virtualSandwich},
\leanid{MPOTensor.TwoSiteStandardFormData.joinedIn\_padBond} and
\leanid{MPOTensor.joinedIn\_standardFormTensorOfGates}. They are assembled in
\leanid{MPOTensor.IsMPU.exists\_joinedIn\_standardFormTensorOfGates} and
\leanid{MPOTensor.IsMPU.equivalent\_of\_index\_eq}. The statement has the
source's hypotheses; the note records only the supplied argument.

\bibliographystyle{alpha}
\bibliography{references}
\end{document}
